%! Author = lili
%! Date = 7/20/26


\newglossaryentry{k01}{type=dinge,
    name={Warum ist \art{A.~thaliana} ein gutes Modellsystem? (8 Punkte)},
    description={Kleine Pflanze ($>$1000 Individuen/m$^2$); \art{Brassicaceae} (Verwandtschaft zu Kohl/Raps); diploid; kleines, wenig repetitives Genom; kurzer Zyklus (6--8 Wochen); bis 10.000 Samen/Pflanze; einfache Mutagenese; Selbstbefruchter mit möglicher Fremdbefruchtung; effektive Transformation}}

\newglossaryentry{k02}{type=dinge,
    name={Was passiert bei der doppelten Befruchtung -- und welche Ploidiegrade entstehen?},
    description={Der Pollenschlauch wächst durch Narbe und Griffel zur Mikropyle und entlässt zwei Spermakerne. Spermakern + Eizelle $\to$ Zygote (2n); Spermakern + 2 Polkerne $\to$ Endospermkern (3n)}}

\newglossaryentry{k03}{type=dinge,
name={Nenne die Reihenfolge der Embryonalstadien.},
description={Zygote $\to$ 2-Zell $\to$ 2/4-Zell-EP $\to$ 8-Zell-EP $\to$ 16-Zell-EP $\to$ Globular $\to$ Transition $\to$ Herz $\to$ Torpedo $\to$ Walking-stick $\to$ reifer Embryo}}

\newglossaryentry{k04}{type=dinge,
name={Welche drei Mutagenesearten gibt es und welchen Mutationstyp erzeugen sie jeweils?},
description={Chemikalien (z.\,B.\ EMS) $\to$ Punktmutationen; Bestrahlung (fast neutrons, Gammastrahlen) $\to$ Deletionen; Transposon- oder T-DNA-Mutagenese $\to$ Insertionen im Genom. Alle drei sind ungerichtet}}

\newglossaryentry{k05}{type=dinge,
name={Warum überhaupt Mutagenese, wenn Mutationen auch spontan entstehen?},
description={Die spontane Mutationsrate ist relativ gering. Rechnung aus der VL: In einem Feld mit 10.000 Pflanzen zu je 1000 Samen finden sich zwar ca.\ 6 Mio.\ Mutationen (5\,\% des Genoms) -- für einen gezielten Screen ist die Ausbeute pro Individuum aber zu niedrig}}

\newglossaryentry{k06}{type=dinge,
name={Unterschied Mutationsrate vs.\ Substitutionsrate?},
description={Mutationsrate = Rate der Genomveränderung durch DNA-Schaden, fehlerhafte Reparatur, Genkonversion, Replikationsfehler. Substitutionsrate = Rate, mit der Mutationen unter Demographie und Selektion überdauern. Unter neutraler Evolution sind beide gleich}}

\newglossaryentry{k07}{type=dinge,
name={Warum sind M1-Pflanzen nach Samenmutagenese Mosaike, F1-Pflanzen nach Pollenmutagenese aber nicht?},
description={Bei der Samenmutagenese wird ein vielzelliger Embryo behandelt -- verschiedene Zellen tragen verschiedene Mutationen. Bei der Pollenmutagenese wird eine einzelne haploide Spermazelle mutagenisiert; alle Zellen der resultierenden F1-Pflanze tragen dieselbe Mutation ($m/+$)}}

\newglossaryentry{k08}{type=dinge,
name={F1 einer Kreuzung Mutante $\times$ Wildtyp zeigt den Mutantenphänotyp. Was folgt daraus?},
description={Die Mutation ist dominant (das Wildtyp-Allel ist rezessiv) -- ein eher seltener Fall. Zeigt die F1 dagegen Wildtyp-Phänotyp, ist die Mutation rezessiv}}

\newglossaryentry{k09}{type=dinge,
name={Kreuzung zweier rezessiver Mutanten X und Y: F1 = Wildtyp. Interpretation?},
description={Komplementation: die Mutationen liegen in unterschiedlichen Genen. F1 = Mutante bedeutet dagegen: keine Komplementation, gleiches Gen, unterschiedliche Allele}}

\newglossaryentry{k10}{type=dinge,
name={Nenne die Stationen auf dem ``langen Weg zum Gen''.},
description={1.~Selektion von Mutanten im Screen (erkennbarer oder selektierbarer Phänotyp); 2.~Kreuzung mit Wildtyp zur Bestimmung des Erbgangs; 3.~genetische Kartierung (Allelie mit bekannten Mutationen; Kopplung mit kartierten Markern $\to$ Chromosom/Bereich); 4.~physische Kartierung (exakte Position, welcher BAC/YAC); 5.~Klonierung und Sequenzierung; 6.~Komplementation}}

\newglossaryentry{k11}{type=dinge,
name={Was ist die Voraussetzung für eine Kartierung mit molekularen Markern?},
description={Eine Population von Individuen, in der genotypische (und ggf.\ die damit verbundenen phänotypischen) Unterschiede aufspalten (segregieren) -- z.\,B.\ die F2 aus Col-0 $\times$ \emph{tt1}/L\emph{er}}}

\newglossaryentry{k12}{type=dinge,
name={Wie unterscheiden sich RFLP und SSLP im Nachweisprinzip?},
description={RFLP: Mutation zerstört/erzeugt eine Restriktionsschnittstelle $\to$ unterschiedliche Fragmentlängen nach Verdau (klassisch mit Southern/Hybridisierung). SSLP: variable Zahl kurzer Repeats $\to$ PCR mit flankierenden Primern liefert unterschiedlich lange Amplifikate; kodominant, Heterozygote zeigen beide Banden}}

\newglossaryentry{k13}{type=dinge,
name={Was ist eine ``unsterbliche'' Kartierungspopulation und wie stellt man sie her?},
description={RILs (rekombinante Inzuchtlinien): eine F2 wird durch ca.\ 7 weitere Selbstungsgenerationen (bis F8) homozygot fixiert. Die Rekombinationsereignisse bleiben dauerhaft erhalten, sodass beliebig viele Labore über Jahre hinweg neue Marker in derselben Population kartieren können}}

\newglossaryentry{k14}{type=dinge,
name={Wie entsteht eine physikalische Karte aus einer BAC-Bank?},
description={Geordnete Klonbank auf Membran repräsentieren $\to$ Sonde vom Ende eines BACs markieren $\to$ hybridisieren $\to$ überlappende BACs identifizieren $\to$ deren Enden als neue Sonden verwenden $\to$ iterieren, bis alle Contigs verknüpft sind $\to$ Assembly}}

\newglossaryentry{k15}{type=dinge,
name={Wie beweist man molekular, dass ein Kandidatengen wirklich das mutierte Gen ist?},
description={Komplementation: Transformation der homozygoten Mutante mit dem Wildtyp-Allel (über Agrobacterium-vermittelten Gentransfer). Zeigen die transgenen Pflanzen Wildtyp-Phänotyp, ist die Genidentität bewiesen}}

\newglossaryentry{k16}{type=dinge,
name={Was wird beim Agrobacterium-vermittelten Gentransfer übertragen -- und wodurch ist es begrenzt?},
description={Die T-DNA, begrenzt durch Left Border und Right Border: zwei fast identische 25-bp-Wiederholungen. Alles dazwischen (inkl.\ eingesetzter Fremdgene und eines Selektionsmarkers) wird übertragen und ins Pflanzengenom integriert}}

\newglossaryentry{k17}{type=dinge,
name={Kenndaten des \art{A.~thaliana}-Genoms (Stand 2000)?},
description={ca.\ 115,4 Mb (1n), 5 Chromosomen (17,5--29,1 Mb); mt 372 kb, cp 153 kb; 10/10/80\,\% hoch/mittel/nicht-niedrig repetitiv; ca.\ 25.500 Gene; Gendichte 1 Gen pro 4,5 kb; Genlänge 2 kb; 5,2 Exons pro Gen; Exon 250 bp, Intron 168 bp}}

\newglossaryentry{k18}{type=dinge,
name={Was bedeutet At3g51240?},
description={AGI-Code: \art{A.~thaliana}, Chromosom 3, gene, laufende Nummer 51240 -- konkret F3H, das \emph{tt6}-Gen}}

\newglossaryentry{k19}{type=dinge,
name={Wie hat sich die Annotation von 2000 bis heute verändert?},
description={Nature 12/2000: 25.498 proteinkodierende Gene, keine Angaben zu alternativem Spleifsen. Araport11 (06/2016): 27.416 proteinkodierende Loci, 10.696 (39\,\%) mit mehreren Spleifsvarianten, deutlich mehr ncRNAs (lincRNA 36 $\to$ 2.444, miRNA 177 $\to$ 325). TAIR12 (06/2026): 26.867 proteinkodierend, 2.008 lncRNA, 558 snoRNA, 526 tRNA, 387 miRNA. Treiber: RNA-seq, Proteomik, manuelle Kuration, Community-Input}}

\newglossaryentry{k20}{type=dinge,
name={Vorwärts- vs.\ Reverse Genetik in einem Satz.},
description={Vorwärts: Mutante/Phänotyp $\rightarrow$ Gen. Reverse: Gen $\rightarrow$ phänotypische Unterschiede}}

\newglossaryentry{k21}{type=dinge,
name={Welche Mutageneseverfahren sind gerichtet?},
description={Homologe Rekombination, Zinkfingernukleasen, TALENs und CRISPR/Cas9 (Genom-Editierung). Ungerichtet sind dagegen Bestrahlung, EMS und Transposon-/T-DNA-Insertion}}

\newglossaryentry{k22}{type=dinge,
name={Wie erzeugt man einen FST?},
description={Genomische DNA schneiden (\emph{Bfa}I/\emph{Csp6}I) $\to$ Linker ligieren $\to$ lineare Amplifikation mit T-DNA-Primer (meist LB) $\to$ PCR mit nested T-DNA-Primer + Linkerprimer $\to$ Sequenzierung mit weiterem nested Primer $\to$ Trace files}}

\newglossaryentry{k23}{type=dinge,
name={Welche bioinformatische Pipeline verarbeitet die FSTs bei GABI-Kat?},
description={Trace files $\to$ Basecalling $\to$ Vector trimming $\to$ Quality trimming $\to$ FASTA $\to$ BLASTN gegen die Annotation (früher TIGR v5, später Araport11) $\to$ BLAST report $\to$ BLAST-Parser $\to$ MySQL-Datenbank mit annotierten FSTs (SimpleSearch)}}

\newglossaryentry{k24}{type=dinge,
name={Wo im Gen inserieren T-DNAs bevorzugt?},
description={Nicht zufällig: Die Insertionsfrequenz ist nahe der Transkriptionsinitiation und nahe der Transkriptionstermination deutlich erhöht (Peak kurz vor dem Startpunkt) gegenüber einer Zufallsverteilung}}

\newglossaryentry{k25}{type=dinge,
name={Drei Beispiele aus der Reverse-Genetik-Praxis (Gen $\to$ Phänotyp)?},
description={At4g18960 (Linie 001G09, Insertion 2.\ Intron) $\to$ ``flower in flower'' $\to$ AGAMOUS. At5g16400 (020E05, 2.\ Intron) $\to$ zwergwüchsig $\to$ Thioredoxin f2. At2g26670 (034B01, 1.\ Exon) $\to$ albino $\to$ Hämoxygenase}}

\newglossaryentry{k26}{type=dinge,
name={Warum reicht ein einzelner FST nicht zur Charakterisierung einer T-DNA-Linie?},
description={Long-read-(ONT-)Sequenzierung zeigte: T-DNA-Insertionen bestehen oft aus komplexen Arrays mehrerer
T-DNAs; im Mittel ca.\ 2 Insertionsloci pro Linie; T-DNA-Mutagenese kann Chromosomenfusionen mit kompensierenden Translokationen und gro{\ss}e Duplikationen auslösen. Für klare Genotyp-Phänotyp-Beziehungen müssen \emph{beide} T-DNA::Genom-Übergänge bestimmt werden}}

\newglossaryentry{k27}{type=dinge,
name={Wie hoch ist der Anteil funktionell charakterisierter \art{Arabidopsis}-Gene?},
description={12/2000: 10\,\% experimentell charakterisiert, 60\,\% funktionelle Kategorie erkennbar, 30\,\% unbekannt. 01/2017: 21\,\% experimentell charakterisiert, 66\,\% Kategorie erkennbar, 13\,\% unbekannt}}
